\documentclass[11pt,a4paper]{article}

\usepackage{fontspec}
\setmainfont{DejaVuSans.ttf}[
  Path=/usr/local/texlive/2024/texmf-dist/fonts/truetype/public/dejavu/,
  BoldFont=DejaVuSans-Bold.ttf,
  ItalicFont=DejaVuSans-Oblique.ttf,
  BoldItalicFont=DejaVuSans-BoldOblique.ttf
]
\usepackage[russian]{babel}
\usepackage{geometry}
\geometry{left=20mm,right=20mm,top=23mm,bottom=20mm}
\usepackage{xcolor}
\usepackage{enumitem}
\usepackage{needspace}
\usepackage{fancyhdr}
\usepackage{hyperref}
\hypersetup{colorlinks=true,linkcolor=black,urlcolor=blue}

\definecolor{navy}{HTML}{183153}
\definecolor{muted}{HTML}{4A5868}
\definecolor{deadline}{HTML}{C62828}
\definecolor{teal}{HTML}{167D84}
\definecolor{pale}{HTML}{EDF3F7}

\pagestyle{fancy}
\fancyhf{}
\fancyhead[L]{\small\bfseries\color{muted}Проект QMIX}
\fancyfoot[R]{\small\color{muted}\thepage}
\renewcommand{\headrulewidth}{0pt}
\setlength{\headheight}{22pt}
\fancypagestyle{firstpage}{\fancyhf{}\fancyfoot[R]{\small\color{muted}\thepage}\renewcommand{\headrulewidth}{0pt}}
\setlength{\parindent}{0pt}
\setlength{\parskip}{2pt}
\setlist[itemize]{leftmargin=1.45em,itemsep=1pt,topsep=3pt,parsep=0pt}

\newcommand{\person}[2]{%
  \Needspace{12\baselineskip}
  \vspace{5pt}
  \noindent\colorbox{pale}{\parbox{\dimexpr\linewidth-2\fboxsep\relax}{%
    \hspace*{7pt}\rule{0pt}{25pt}{\Large\bfseries\color{navy}#1}%
    \hfill{\bfseries\color{deadline}Дедлайн: #2}\hspace*{7pt}}}\par
  \vspace{1pt}
}

\begin{document}

\thispagestyle{firstpage}
{\Huge\bfseries\color{navy}Проект QMIX}\par
\vspace{3pt}
{\color{teal}\rule{44mm}{2.2pt}}\par
\vspace{3pt}

\person{Алина}{10 октября}
\begin{itemize}
  \item Подготовить 4--5 слайдов proposal presentation: какую задачу решает QMIX; что вычисляет каждый агент $Q_i$; как mixing network получает $Q_{\mathrm{tot}}$; зачем при обучении нужно полное состояние и что означает условие монотонности.
  \item На отдельном слайде описать Simple Spread: цель агентов, их наблюдения, пять возможных действий и общую награду. Параметры среды сверить с результатом Жени.
  \item Нарисовать схему «наблюдения $\rightarrow Q_i \rightarrow$ mixing network $\rightarrow Q_{\mathrm{tot}}$» и отметить, какие части используются при обучении, а какие — при выполнении политики.
\end{itemize}
\textbf{Сдать:} ссылку на редактируемую proposal presentation и готовую схему QMIX.\par
\textbf{Кому передать:} Дане — слайды и схему для финальной презентации.\par

\person{Женя}{8 октября}
\begin{itemize}
  \item Запустить Simple Spread в Google Colab: два агента, дискретные действия, параметр \texttt{local\_ratio=0} для общей награды.
  \item Вывести и подписать пример наблюдения каждого агента, доступные действия, полное состояние и награду после одного шага. Проверить, что награда у обоих агентов одинаковая.
  \item Провести один полный эпизод со случайными действиями и сохранить изображение или короткую запись эпизода.
  \item Записать точные параметры среды и версии библиотек, с которыми запуск работает.
\end{itemize}
\textbf{Сдать:} ссылку на запускаемый Colab, визуальный пример и краткое описание параметров среды.\par
\textbf{Кому передать:} Кариму — рабочие ячейки и параметры; Алине — описание среды для proposal presentation.\par

\person{Карим}{10 октября}
\begin{itemize}
  \item Собрать основной Colab на основе среды Жени. В notebook должны запускаться обучение QMIX с нуля, сохранение модели и построение графика награды в ходе обучения.
  \item Показать, где вычисляются индивидуальные $Q_i$, как они объединяются в $Q_{\mathrm{tot}}$, как считается TD-цель и какая ошибка минимизируется.
  \item Сделать один полностью числовой пример TD-обновления: задать награду, $\gamma$, текущую и следующую оценку; вычислить TD-цель, ошибку и одно изменение оценки или параметра. Показать промежуточные числа.
  \item Загрузить сохранённую модель и провести с ней эпизод без дополнительного обучения.
\end{itemize}
\textbf{Сдать:} ссылку на основной Colab, checkpoint обученной модели и график обучения.\par
\textbf{Кому передать:} Аслану — checkpoint и код запуска модели; Дане — notebook и график для презентации.\par

\person{Аслан}{11 октября}
\begin{itemize}
  \item Оценить checkpoint Карима без случайного выбора действий на фиксированном наборе, например на 50 эпизодах.
  \item На тех же начальных состояниях запустить случайных агентов и сравнить результаты.
  \item Посчитать среднюю награду и разброс для обоих вариантов. Если позволяет время, повторить обучение QMIX с 2--3 разными seed; фактическое число запусков указать в результатах.
  \item Построить график сравнения, записать один эпизод обученных агентов для видео и указать, какой checkpoint использован.
\end{itemize}
\textbf{Сдать:} таблицу чисел, график, запись эпизода и краткий вывод на 2--3 предложения.\par
\textbf{Кому передать:} Дане — результаты и запись для финальной презентации; Кариму — числа для notebook.\par

\person{Даня}{12 октября — материалы; 13 октября — проверка}
\begin{itemize}
  \item Собрать финальные слайды из материалов команды: краткое введение, QMIX, Simple Spread, схема их взаимодействия, числовой пример Карима, реальные результаты Аслана, демонстрация, выводы и источники.
  \item Подготовить сценарий на 5--7 минут и записать видео с рассказом по слайдам и показом эпизода обученных агентов.
  \item Собрать ссылки на предварительную презентацию, основной Colab, checkpoint, финальные слайды и видео. Закончить комплект до конца 12 октября.
  \item 13 октября проверить открытие ссылок, звук и длительность видео, а также запуск notebook в чистой сессии Colab. Исправить найденные ошибки; 14 октября отправить комплект до дедлайна.
\end{itemize}
\textbf{Сдать:} финальную презентацию, видеофайл на 5--7 минут и проверенный комплект ссылок.\par
\textbf{Кому передать:} всей команде — для последней проверки; затем преподавателю через установленный способ сдачи.\par

\end{document}
